\documentclass[a4paper,12pt]{article}
\usepackage[T2A]{fontenc}
\usepackage[utf8]{inputenc}
\usepackage[russian]{babel}
\usepackage{amsmath,amssymb,mathtools}
\usepackage{helvet}
\renewcommand{\familydefault}{\sfdefault}
\usepackage{geometry}
\geometry{left=24mm,right=24mm,top=25mm,bottom=24mm}
\usepackage{microtype}
\usepackage{icomma}
\usepackage{tikz}
\usetikzlibrary{arrows.meta,positioning,shapes.geometric,calc}
\usepackage{booktabs,tabularx,array}
\usepackage{enumitem}
\usepackage{setspace}
\usepackage{xcolor}
\usepackage{fancyhdr}
\usepackage{hyperref}
\hypersetup{hidelinks,
  pdftitle={Чек-лист: прямая пропорциональность и построение графиков},
  pdfauthor={Лёвин Артём Александрович}}
\definecolor{accent}{HTML}{008080}
\definecolor{darkaccent}{HTML}{005757}
\definecolor{textgray}{HTML}{30373A}
\definecolor{muted}{HTML}{637073}
\color{textgray}
\setstretch{1.20}
\setlength{\parindent}{0pt}
\setlength{\parskip}{0.55em}
\setlength{\headheight}{14pt}
\pagestyle{fancy}
\fancyhf{}
\fancyhead[L]{\footnotesize\color{muted} Прямая пропорциональность}
\fancyhead[R]{\footnotesize\color{muted} Кирилл Зиновьев}
\fancyfoot[C]{\footnotesize\color{muted}\thepage}
\renewcommand{\headrulewidth}{0.35pt}
\renewcommand{\footrulewidth}{0pt}
\newcommand{\sect}[1]{\vspace{0.55em}\section*{\color{darkaccent}\large #1}\vspace{-0.6em}{\color{accent}\rule{\linewidth}{0.45pt}}\vspace{0.3em}}
\newcommand{\subt}[1]{\vspace{0.45em}{\color{darkaccent}\bfseries #1}\par\vspace{-0.2em}}
\newcommand{\idea}[1]{{\color{darkaccent}\bfseries Ключевая идея.} #1}
\newcommand{\task}[1]{\item #1\vspace{0.35em}}
\setlist[enumerate]{leftmargin=*,itemsep=0.3em,topsep=0.4em}

\begin{document}
% Титульный лист
\thispagestyle{empty}
\begin{center}
\vspace*{33mm}
{\small\color{darkaccent}\MakeUppercase{Математика}}\par
\vspace{16mm}
{\LARGE Чек-лист}\par
\vspace{12mm}
{\Large\color{accent} Прямая пропорциональность\\[3pt] и построение графиков}\par
\vspace{12mm}
{\normalsize 7 класс. Алгебра}\par
\vfill
{\large Кирилл Зиновьев}\par
\vspace{10mm}
{4 ноября 2026 г.}\par
\vspace{21mm}
{\small Лёвин Артём Александрович\\ эксклюзивно для Кирилла Зиновьева}\par
\vspace{18mm}
\begin{tikzpicture}[x=1cm,y=1cm]
\draw[accent,line width=0.7pt] (0,0)
 .. controls (1.1,.36) and (2.0,-.26) .. (3.2,0)
 .. controls (4.4,.28) and (5.3,-.28) .. (6.5,0)
 .. controls (7.7,.3) and (8.7,-.33) .. (10,0);
\end{tikzpicture}
\end{center}
\clearpage

\sect{1. Что такое прямая пропорциональность}
Прямая пропорциональность задаётся формулой
\[
 y=kx,
\]
где $x$ --- независимая переменная, $y$ --- соответствующее ей значение,
$k$ --- постоянный коэффициент пропорциональности ($k\ne 0$ при построении обычной наклонной прямой).

График --- прямая, проходящая через начало координат $O(0;\,0)$.
Поэтому для её построения достаточно двух различных точек.

\idea{Значения $x$ можно выбрать самостоятельно. Значения $y$ вычисляют по заданной формуле.}

\sect{2. Таблица значений: как её заполнить}
Пусть дана формула $y=3x$. Удобно взять $x=0$ и $x=1$:
\begin{center}
\renewcommand{\arraystretch}{1.28}
\begin{tabular}{@{}lcc@{}}
\toprule
$x$ & $0$ & $1$\\
\midrule
$y=3x$ & $3\cdot0=0$ & $3\cdot1=3$\\
\bottomrule
\end{tabular}
\end{center}
Получаем точки $A(0;\,0)$ и $B(1;\,3)$. На координатной плоскости сначала
отмечают $A$ и $B$, затем проводят прямую через обе точки.

\subt{Почему нуль удобен}
Если $x=0$, то при любой формуле $y=kx$ имеем $y=k\cdot0=0$. Первая точка всегда $O(0;\,0)$.

\subt{Как читать координаты}
Запись $B(2;\,-3)$ означает: от начала координат вправо на $2$ единицы по оси $x$,
затем вниз на $3$ единицы по оси $y$.
\clearpage

\sect{3. Отрицательный коэффициент и знаки}
Для $y=-1{,}5x$ выбираем $x=0$ и $x=2$:
\[
\begin{array}{c|cc}
 x&0&2\\\hline
 y=-1{,}5x&0&-3
\end{array}
\qquad\Longrightarrow\qquad A(0;\,0),\ B(2;\,-3).
\]
Когда $x>0$, получаем $y<0$. Когда $x<0$, получаем $y>0$, поскольку минус на минус даёт плюс.
\[
\text{Например:}\quad x=-2\quad\Longrightarrow\quad y=-1{,}5\cdot(-2)=3.
\]

\sect{4. Как выглядит полученная прямая}
\begin{center}
\begin{tikzpicture}[scale=0.87,>=Stealth]
  \draw[->,line width=0.9pt] (-4.1,0)--(4.1,0) node[right] {$x$};
  \draw[->,line width=0.9pt] (0,-4.6)--(0,4.5) node[above] {$y$};
  \foreach \x in {-3,-2,-1,1,2,3} {
    \draw (\x,.07)--(\x,-.07);
    \node[below,font=\footnotesize] at (\x,-.12) {$\x$};
  }
  \foreach \y in {-4,-3,-2,-1,1,2,3,4} {
    \draw (.07,\y)--(-.07,\y);
    \node[left,font=\footnotesize] at (-.12,\y) {$\y$};
  }
  \draw[accent,very thick,<->] (-2.8,4.2)--(2.8,-4.2);
  \fill[darkaccent] (0,0) circle (2.6pt);
  \node[above right,font=\small] at (0,0) {$A(0;\,0)$};
  \fill[darkaccent] (2,-3) circle (2.6pt);
  \node[right,font=\small] at (2,-3) {$B(2;\,-3)$};
\end{tikzpicture}
\end{center}
\textbf{Важно:} точки $A$ и $B$ помогают провести график. Сама прямая продолжается
за эти точки в обе стороны. Отрезок $AB$ не является полным графиком функции.
\clearpage

\sect{5. Что проверять в своей работе}
\begin{enumerate}[label=\textbf{\arabic*.}]
\item В таблице подписаны строки $x$ и $y$, даже если в формуле есть минус.
\item Выбраны два разных значения $x$.
\item Каждое значение $y$ найдено подстановкой выбранного $x$ в формулу.
\item При отрицательных $x$ правильно учтено правило знаков.
\item Точки отмечены в соответствии с парами $(x;\,y)$.
\item Через точки проведена прямая, продолженная в обе стороны.
\end{enumerate}

\sect{6. Ошибки, на которые стоит обратить внимание}
\renewcommand{\arraystretch}{1.36}
\begin{tabularx}{\linewidth}{@{}p{0.39\linewidth}X@{}}
\toprule
Ошибочная запись или действие & Как исправить\\
\midrule
В заголовке второй строки таблицы пишут $-x$ вместо $y$. & Строки всегда подписывают $x$ и $y$. Знак учитывается при вычислении $y$.\\
При $y=-x$ и $x=-2$ получают $y=-2$. & Подставьте $x=-2$: $y=-(-2)=2$.\\
Две точки соединили отрезком. & Продлите линию за обе точки: график --- прямая.\\
\bottomrule
\end{tabularx}
\clearpage

\begin{center}
{\Large\color{darkaccent} Алгоритм: построение графика $y=kx$}\par
\vspace{9mm}
\end{center}
\begin{center}
\begin{tikzpicture}[
  node distance=11mm,
  start/.style={ellipse,draw=darkaccent,very thick,align=center,minimum width=47mm,minimum height=10mm,font=\small},
  action/.style={rectangle,rounded corners=2pt,draw=accent,thick,align=center,text width=91mm,minimum height=13mm,font=\small,inner sep=4pt},
  check/.style={diamond,aspect=2.1,draw=darkaccent,thick,align=center,text width=39mm,inner sep=1.5pt,font=\small},
  arr/.style={-{Stealth[length=2.5mm]},thick,draw=darkaccent}
]
  \node[start] (begin) {Дана формула $y=kx$};
  \node[action,below=of begin] (choose) {Составьте таблицу: выберите два значения $x$.\\Удобно одно из них взять равным нулю};
  \node[check,below=12mm of choose] (check) {Выбраны разные значения $x$?};
  \node[action,below=17mm of check] (compute) {Подставьте каждое $x$ в формулу,\\ вычислите оба значения $y$};
  \node[action,below=of compute] (plot) {Отметьте на плоскости две точки\\$(x_1;\,y_1)$ и $(x_2;\,y_2)$};
  \node[action,below=of plot] (line) {Проведите через точки прямую\\и продолжите её в обе стороны};
  \node[start,below=of line] (finish) {График готов};
  \draw[arr] (begin)--(choose);
  \draw[arr] (choose)--(check);
  \draw[arr] (check)--node[right,font=\footnotesize]{да}(compute);
  \draw[arr] (compute)--(plot);
  \draw[arr] (plot)--(line);
  \draw[arr] (line)--(finish);
  \coordinate (goLeft) at ($(check.west)+(-21mm,0)$);
  \coordinate (toChoose) at ($(choose.west)+(-21mm,0)$);
  \draw[arr] (check.west)--node[above,font=\footnotesize]{нет}(goLeft) |- (toChoose) -- (choose.west);
\end{tikzpicture}
\end{center}
\clearpage

\sect{Тренировочный блок. Путь А}
Решайте последовательно. В заданиях на построение сначала составляйте таблицу,
затем отмечайте точки и проводите полную прямую. Используйте тетрадь в клетку.

\subt{Базовый уровень}
\begin{enumerate}[label=\textbf{\arabic*.}]
\task{Заполните строку $y$ для формулы $y=3x$ при $x=-1,\ 0,\ 2$.}
\task{Заполните строку $y$ для формулы $y=-x$ при $x=-3,\ 0,\ 4$.}
\task{Заполните строку $y$ для формулы $y=-2x$ при $x=-2,\ 0,\ 1$.}
\subt{Средний уровень}
\task{Постройте график $y=4x$ по точкам с абсциссами $x=0$ и $x=1$. Запишите координаты точек.}
\task{Постройте график $y=1{,}5x$ по точкам с абсциссами $x=0$ и $x=2$. Запишите координаты.}
\task{Постройте график $y=-1{,}5x$ по точкам с абсциссами $x=-2$ и $x=0$. Запишите координаты.}
\task{Постройте график $y=2{,}5x$ по точкам с абсциссами $x=0$ и $x=2$. Запишите координаты.}
\end{enumerate}
\clearpage

\sect{Уровень выше среднего}
\begin{enumerate}[label=\textbf{\arabic*.},start=8]
\task{Постройте график $y=-0{,}5x$ по точкам с абсциссами $x=-4$ и $x=0$. Запишите координаты.}
\task{Для $y=-3x$ отмечены $A(0;\,0)$ и $B(2;\,-6)$. Ученик соединил их только отрезком $AB$. Какую ошибку он допустил при построении графика? Вычислите также $y$, если $x=-3$.}
\task{На одной координатной плоскости постройте прямые $y=2x$ и $y=-2x$. Для каждой используйте $x=-2$, $x=0$, $x=2$. Запишите координаты точек и сравните, куда идут прямые при движении вправо.}
\end{enumerate}
\vspace{1em}
\textbf{Проверка после построения.} Для каждой прямой $y=kx$ убедитесь, что она проходит через $O(0;\,0)$ и продолжается за обе отмеченные точки.
\clearpage

\sect{Ответы к тренировочному блоку}
\renewcommand{\arraystretch}{1.26}
\begin{tabularx}{\linewidth}{@{}>{\bfseries}p{12mm}X@{}}
\toprule
№ & Ответ или контрольные точки\\
\midrule
1 & $y=-3,\ 0,\ 6$ (для $x=-1,\ 0,\ 2$).\\
2 & $y=3,\ 0,\ -4$ (для $x=-3,\ 0,\ 4$).\\
3 & $y=4,\ 0,\ -2$ (для $x=-2,\ 0,\ 1$).\\
4 & $(0;\,0)$ и $(1;\,4)$.\\
5 & $(0;\,0)$ и $(2;\,3)$.\\
6 & $(-2;\,3)$ и $(0;\,0)$.\\
7 & $(0;\,0)$ и $(2;\,5)$.\\
8 & $(-4;\,2)$ и $(0;\,0)$.\\
9 & Нужно продлить прямую за обе точки; при $x=-3$ имеем $y=9$.\\
10 & $y=2x$: $(-2;\,-4)$, $(0;\,0)$, $(2;\,4)$; прямая идёт вверх вправо. Для $y=-2x$: $(-2;\,4)$, $(0;\,0)$, $(2;\,-4)$; прямая идёт вниз вправо.\\
\bottomrule
\end{tabularx}

\subt{Короткая памятка}
\renewcommand{\arraystretch}{1.20}
\begin{tabularx}{\linewidth}{@{}p{0.32\linewidth}X@{}}
\toprule
Ситуация & Действие\\
\midrule
Дана формула $y=kx$ & Выберите два разных $x$, вычислите $y$.\\
$x=0$ & Сразу получаете точку $(0;\,0)$.\\
$k<0$ & При положительном $x$ значение $y$ отрицательно.\\
Точки поставлены & Проведите через них прямую в обе стороны.\\
\bottomrule
\end{tabularx}
\vfill
\begin{center}
\color{muted}\footnotesize Лёвин Артём Александрович\\ эксклюзивно для Кирилла Зиновьева
\end{center}
\end{document}
